%%%% FromHawking_DLIndaba2026.tex
\typeout{From Hawking to Intelligent Markets -- DL Indaba 2026}

\documentclass{article}
\pdfpagewidth=8.5in
\pdfpageheight=11in

\usepackage{ijcai26}
\usepackage{times}
\usepackage{soul}
\usepackage{url}
\usepackage[hidelinks]{hyperref}
\usepackage[utf8]{inputenc}
\usepackage[small]{caption}
\usepackage{graphicx}
\usepackage{amsmath}
\usepackage{amsthm}
\usepackage{booktabs}
\usepackage{algorithm}
\usepackage{algorithmic}
\usepackage{multirow}
\usepackage{enumitem}
\usepackage{xcolor}
\usepackage[switch]{lineno}

% Comment out for camera-ready
\linenumbers

% Enable subsubsection numbering (ijcai26.sty sets this to 2)
\setcounter{secnumdepth}{3}

\urlstyle{same}

\newtheorem{definition}{Definition}
\newtheorem{proposition}{Proposition}

\pdfinfo{
/TemplateVersion (IJCAI.2026.0)
}

\title{Privacy-Preserving Market Analytics for Informal Urban Economies in Africa}

\author{
Anonymous Authors
}

\begin{document}

\maketitle

\begin{abstract}
As African cities pursue smart city agendas, modern market initiatives seek to relocate street vendors into standardised hubs. While these policies target improved sanitation and urban order, they ignore the economic logic of subsistence-level retailers, primarily women and youth who rely on the high foot traffic generated by street hawking.
This misalignment creates an inclusion gap, with high entry barriers emerging from exorbitant rent fees, taxes, licenses, and reduced customer access, increasing the risk of vendor displacement a descent into poverty. This paper presents a dual-intervention framework.
First, we develop a federated learning system for distributed market intelligence that preserves vendor data sovereignty while also enabling collective demand--supply insights.
Second, we introduce AI-optimised modular retail units that address the economic and spatial barriers of traditional market infrastructure.
We formalise the \emph{inclusion gap} as the per-vendor economic penalty of lacking formal infrastructure, and derive digital demand-pull and risk-weighted allocation models.
Experiments on Rwanda's Establishment Census of 269,326 enterprises, where 87.8\% of these are informal, World Bank Enterprise Surveys, WFP market prices across Rwanda and Uganda, WorldPop density grids, and OpenStreetMap data indicate that:
(i)~initial experiments suggested high separability; after leakage controls, realistic formalization performance remained competitive but substantially lower;
(ii)~scenario analysis suggests potential income gains under assumed occupancy and demand conditions;
(iii)~risk-weighted spatial allocation improves on population-proportional placement under identical budgets.
A comparative Rwanda--Uganda benchmark on 629 market clients provides evidence of cross-border generalisability.
\end{abstract}

\noindent Keywords: Federated Learning, Market Intelligence, Informal Economies, Digital Visibility, Modular Retail Infrastructure, Inclusive Urban Development.

%=================================================================
\section{Introduction}
\label{sec:intro}
%=================================================================

Across sub-Saharan Africa, informal enterprises constitute the dominant economic form.
In Rwanda, the 2023 Establishment Census records 269,326 enterprises, of which 229,764 (87.8\%) operate informally \cite{nisr2023ec}.
A striking 92.2\% are micro-scale (1--3 workers), and 63.2\% hold capital below 300,000~RWF (\textasciitilde\$214).
Municipal smart city strategies, such as the Kigali Master Plan 2050 \cite{kigali2019master}, prescribe relocation into modern market halls with formalised rent, licensing, and sanitation standards.
However, annual rent in traditional modern markets reaches 1,200,000~RWF, compared with median informal turnover of only 250,000~RWF/year.
This rent-to-revenue mismatch makes formal markets \emph{economically unviable} for the overwhelming majority of informal vendors, creating what we term the \emph{inclusion gap}.

This paper makes three contributions:
\begin{enumerate}[label=\alph*)]
    \item We introduce a federated learning benchmark for informal market analytics using real African datasets (Rwanda Establishment Census, WFP prices, WBES).
    \item We demonstrate that privacy-preserving federated collaboration improves cross-market price forecasting over isolated local models.
    \item We propose an equity-aware decision framework for modular retail unit deployment, evaluated through spatial optimization on district-level data.
\end{enumerate}

%=================================================================
\section{Related Work}
\label{sec:background}
%=================================================================

\subsection{Urban Informality and Forced Formalisation}

Street vending is the primary livelihood for over 2 billion workers globally, with informal employment accounting for 55--80\% of urban jobs in African cities \cite{ilo2023weso}.
Chen~\shortcite{chen2023rethinking} argues that informality is not a transitional phase but a structural feature of developing economies that requires integration rather than elimination.
Forced relocation policies, often termed ``clean-up'' or ``decongestion'' campaigns, have been documented across Johannesburg, Nairobi, Accra, and Kigali \cite{roever2023street}.
These interventions typically reduce vendor income by removing access to high-traffic locations while imposing unaffordable formal costs.
The COVID-19 pandemic further exposed the vulnerability of informal workers \cite{skinner2021covid}, with Adhikari and Shrestha~\shortcite{adhikari2023vendors} finding that post-pandemic recovery has been slowest for vendors displaced into formal markets.
Herrera and L\'{e}autier~\shortcite{herrera2023informal} demonstrate through spatial analysis that informality correlates strongly with urbanisation rate and infrastructure deficits across sub-Saharan Africa.

\subsection{Federated Learning for Resource-Constrained Settings}

Federated learning (FL) enables model training across distributed clients without centralising raw data \cite{mcmahan2017communication,kairouz2021advances}.
FedAvg aggregates locally trained model parameters via weighted averaging, while FedProx \cite{li2020federated} adds a proximal regularisation term to handle statistical heterogeneity.
Recent comprehensive surveys \cite{zhang2022fedsurvey,wen2023survey} identify data heterogeneity, communication efficiency, and privacy as key FL challenges, all of which are acute in informal market settings where vendors are privacy-sensitive and connectivity is intermittent.
Pfitzner et al.~\shortcite{pfitzner2021fedhealth} demonstrate FL's promise in healthcare settings with analogous privacy constraints, while Nguyen et al.~\shortcite{nguyen2022fliot} survey FL for resource-constrained IoT devices.
However, applications of FL to informal market systems in developing economies remain entirely unexplored.

\subsection{Digital Platforms and Spatial Market Design}

Digital platforms have shown promise in connecting informal micro-enterprises to broader markets.
Mobile money, which has reached over 1.75 billion registered accounts in sub-Saharan Africa \cite{suri2023fintech}, provides a ready digital layer for vendor transactions.
However, most digital platforms require centralised data collection that excludes privacy-conscious vendors.
Spatial allocation of market infrastructure has traditionally relied on population-proportional heuristics or maximal covering location models.
Osei et al.~\shortcite{osei2022smartafrica} find that African smart city initiatives often neglect informal economic actors in spatial planning.
Our work departs by deriving viability-aware allocation scores from a mathematical framework that integrates demand density, digital readiness, and infrastructure cost.
Fox~\shortcite{fox2022urbanization} highlights the need for development approaches that account for the political economy of urbanisation. Our modular unit design responds to this call by offering a low-cost, incrementally deployable alternative to mega-market projects.

%=================================================================
\section{Methods}
\label{sec:methodology}
%=================================================================

\subsection{The Inclusion Gap}
\label{sec:problem}

We begin by formalising the economic concepts at the core of our framework, making each metric intuitive before stating its mathematical form.

\subsubsection{Vendor Viability}
For a vendor $v$ operating under scenario $s \in \{\textit{street}, \textit{modular}, \textit{modern}\}$, the net annual viability is:
\begin{equation}
    N(v,s) = \rho(s) \cdot \bar{m} \cdot d - C_s
    \label{eq:viability}
\end{equation}
where $\rho(s)$ is the scenario-specific revenue, $\bar{m}$ is the average margin rate, $d$ is trading days per year, and $C_s$ is the annual cost.
Using parameters from Rwanda's EC 2023 (Table~\ref{tab:viability}), street hawking yields $N_{\textit{street}} = 25{,}000$~RWF/year, while AI-optimised modular units yield $N_{\textit{modular}} = 177{,}500$~RWF/year.
Traditional modern markets are \emph{unviable}: $N_{\textit{modern}} = -720{,}000$~RWF/year.

\subsubsection{District-Level Inclusion Gap ($\Gamma_d$)}
$\Gamma_d$ quantifies the annual economic penalty (in RWF) that each informal vendor in district $d$ suffers from lacking access to formal market infrastructure:
\begin{equation}
    \Gamma_d = \frac{N_d^{\textit{unserved}}}{N_d^{\textit{inf}}} \cdot \max\!\bigl(0,\; N_{\textit{modular}} - N_{\textit{street}}\bigr)
    \label{eq:gamma}
\end{equation}
Here, $N_d^{\textit{inf}}$ is the total number of informal enterprises in district $d$, and $N_d^{\textit{unserved}} = N_d^{\textit{inf}} - N_d^{\textit{served}}$ counts those without any formal or modular infrastructure.
The term $N_{\textit{modular}} - N_{\textit{street}}$ is the modeled income gain associated with switching from street hawking (25,000~RWF/yr) to a modular unit (177,500~RWF/yr) under the stated assumptions.
In other words, the ratio $N_d^{\textit{unserved}} / N_d^{\textit{inf}}$ is the fraction of vendors being ``left behind,'' and the second term is the per-vendor income they are missing out on.
The national average is $\Gamma = 149{,}175$~RWF/year per vendor, totalling 33,820~M~RWF (\textasciitilde\$24.2M) in annual lost income across Rwanda.

\subsubsection{Digital Demand-Pull Coefficient ($\beta_d$)}
$\beta_d$ captures how much digital infrastructure in district $d$ boosts vendor revenue through better demand matching.
We first compute a digital readiness index:
\begin{equation}
    D(d) = 0.3\,M_d + 0.2\,I_d + 0.3\,L_d + 0.2\,(1 - E_d)
    \label{eq:digital}
\end{equation}
where $M_d$ is mobile money penetration, $I_d$ is internet access rate, $L_d$ is digital literacy rate, and $E_d$ is financial exclusion rate, which is inverted so that lower exclusion means higher readiness (all from FinScope \cite{finscope2020rwanda}).
The demand-pull coefficient is then:
\begin{equation}
    \beta_d = 1 + \delta \cdot D(d)
    \label{eq:beta}
\end{equation}
with $\delta = 0.85$ calibrated from Kigali--Kampala market data.
When $D(d) = 0$ (no digital readiness), $\beta_d = 1$ and there is no revenue boost.
At full readiness ($D(d) = 1$), $\beta_d$ reaches 1.85.
In practice, digitally ready urban districts such as Gasabo or Nyarugenge have higher $\beta_d$, meaning vendors there earn more because mobile money payments, online visibility, and demand matching amplify foot traffic.
This is why the digital divide \emph{widens} the inclusion gap: districts with low $\beta_d$ benefit less from modular units.

\subsubsection{Viability Risk Score ($R(d)$)}
$R(d)$ estimates the probability that a modular retail unit placed in district $d$ will fail economically:
\begin{equation}
    R(d) = \frac{\max\!\bigl(0,\; C_{\textit{mod}} - \beta_d \cdot \text{Rev} \cdot \bar{m}\bigr)}{C_{\textit{mod}}} \in [0,1]
    \label{eq:risk}
\end{equation}
where $C_{\textit{mod}}$ is the annual cost of operating a modular unit, $\text{Rev}$ is expected revenue, and $\bar{m}$ is the average profit margin.
If digitally-boosted revenue covers costs ($\beta_d \cdot \text{Rev} \cdot \bar{m} \geq C_{\textit{mod}}$), then $R(d) = 0$, meaning there is no risk.
If revenue falls short, $R(d)$ increases toward~1, proportional to the shortfall.
The allocation score $S(d) = N_d^{\textit{inf}} \times (1 - R(d)) \times \beta_d$ then determines how many modular units each district receives: districts with high risk get fewer units, reducing allocation to areas where modular units are less likely to be economically viable.

\subsection{Data Sources}

We integrate six data layers spanning two countries (Table~\ref{tab:data}):

\begin{table}[t]
    \centering
    \caption{Multi-source data stack.}
    \label{tab:data}
    \begin{tabular}{@{}llrl@{}}
        \toprule
        Source & Country & Records & Role \\
        \midrule
        NISR EC 2023 & Rwanda & 269,326 & Inclusion gap \\
        WBES 2023 & Rwanda & 358 & Enterprise features \\
        WFP Prices & Rwanda & 100,733 & Price time series \\
        WFP Prices & Uganda & 25,526 & Comparative TS \\
        WorldPop & Both & 270K cells & Demand proxy \\
        OSM & Both & 598 mkts & Market graph \\
        FinScope & Both & 46 districts & Digital readiness \\
        \bottomrule
    \end{tabular}
\end{table}

Enterprise features are drawn from the WBES Rwanda 2023 dataset \cite{worldbank2023wbes} (358 formal enterprises, 355 variables), from which we extract 87 variables across nine thematic groups: registration, infrastructure, finance, labour, sales, innovation, performance, obstacles, and identity.
After removing columns with $>$50\% missingness, 63 numeric features remain.
Sentinel values ($-9$, $-8$, $-7$) are mapped to NaN and imputed with column medians.

For market price data, WFP VAM \cite{wfp2024vam} provides 100,733 monthly commodity price observations across 109 Rwandan markets (2000--2026) and 25,526 observations across 43 Ugandan markets (2006--2026).

The spatial layers combine WorldPop \cite{tatem2017worldpop} 1\,km$^2$ population density grids (25,826 cells for Rwanda, 243,717 for Uganda), which serve as catchment demand proxies.
OpenStreetMap \cite{haklay2008openstreetmap} provides 50 market nodes and 569 primary road segments for Rwanda, and 548 markets with 1,135 road segments for Uganda.

\subsection{Federated Learning Framework}

\subsubsection{Task 1: Formalization Prediction}

For formalization prediction, MarketNet is a 3-layer feedforward network: $\text{input} \to 64 \to 32 \to 1$ with ReLU activations, dropout ($p{=}0.3$), and BCEWithLogitsLoss.

Enterprises are assigned to $K{=}30$ simulated district-level clients, matching Rwanda's administrative districts where sufficient rows are available and using stratified pseudo-district partitions otherwise.
All primary metrics are reported over five repeated stratified random splits.

The binary formalization label is constructed from WBES variable \texttt{b3a} (registered when starting operations): class~1 (formal, $n{=}294$) and class~0 (informal, $n{=}64$).
To avoid label leakage, all registration-status variables and direct proxies were excluded from predictive features.
Results reported reflect leakage-controlled evaluation.

We compare four training paradigms:
(i)~\emph{Centralized}: all data pooled, $T{=}150$ epochs;
(ii)~\emph{Local-only}: each client trains independently;
(iii)~\emph{FedAvg}: 30 rounds $\times$ 5 local epochs, LR~$= 0.005$;
(iv)~\emph{FedProx}: as FedAvg with proximal term $\mu{=}0.1$.

To address the small sample size (358 enterprises), we evaluate feature selection and SMOTE \cite{chawla2002smote} only as sensitivity checks.
Primary formalization results use the original WBES rows after leakage controls; SMOTE-expanded experiments are reported as exploratory because synthetic samples may inflate apparent generalisation.

\subsubsection{Task 2: Federated Price Prediction}

For price prediction, PriceNet is a regression network: $6 \to 32 \to 16 \to 1$ with ReLU activations and MSELoss.
Input features are sliding windows of 6 monthly prices.

Each WFP market$\times$commodity pair constitutes a federated client.
We select the top-5 commodities by observation count, yielding 426 clients for Rwanda (14,046 train / 3,729 test samples) and 203 clients for Uganda (7,986 train / 2,102 test).

The price-forecasting evaluation uses a strict temporal train/test partition in which the first 80\% of months form the training set and the last 20\% form the held-out test set, preventing information leakage.

Each market node is enriched with WorldPop catchment population within a 10\,km radius, OSM retail density, and FinScope digital readiness score.
Spatial graphs connect markets within 80\,km, producing graphs with average degrees of 169.1 (Rwanda) and 12.4 (Uganda).

\subsection{Spatial Allocation of Modular Units}

Given a budget of $B{=}100$ modular units, each serving $C{=}50$ vendors, we allocate units to districts $d \in \mathcal{D}$ proportional to the allocation score:
\begin{equation}
    u_d = \left\lfloor B \cdot \frac{S(d)}{\textstyle\sum_{d'} S(d')} \right\rfloor
    \label{eq:allocation}
\end{equation}
Remaining units (from rounding) are greedily assigned to districts with highest residual score.

\begin{algorithm}[t]
\caption{Federated Market Intelligence Training}
\label{alg:fedmarket}
\begin{algorithmic}[1]
\REQUIRE $K$ clients, $T$ rounds, local epochs $E$, learning rate $\eta$, proximal weight $\mu$
\ENSURE Global model parameters $\theta^T$
\STATE Initialise global model $\theta^0$
\FOR{$t = 0$ \TO $T-1$}
    \FOR{each client $k \in \{1, \ldots, K\}$ in parallel}
        \STATE $\theta_k \leftarrow \theta^t$
        \FOR{epoch $e = 1$ \TO $E$}
            \FOR{batch $(x, y) \in \mathcal{D}_k$}
                \STATE $\mathcal{L} \leftarrow \ell(\theta_k; x, y) + \frac{\mu}{2}\|\theta_k - \theta^t\|^2$
                \STATE $\theta_k \leftarrow \theta_k - \eta \nabla \mathcal{L}$
            \ENDFOR
        \ENDFOR
    \ENDFOR
    \STATE $\theta^{t+1} \leftarrow \sum_{k=1}^{K} \frac{|\mathcal{D}_k|}{|\mathcal{D}|} \theta_k$ \COMMENT{Weighted aggregation}
\ENDFOR
\end{algorithmic}
\end{algorithm}

%=================================================================
\section{Experiments and Results}
\label{sec:results}
%=================================================================

\subsection{Formalization Prediction}

Table~\ref{tab:formal} reports formalization prediction results.
Initial experiments suggested high separability; after leakage controls, realistic performance remained competitive but substantially lower.
FedAvg and FedProx are comparable to centralized neural baselines and are most useful as privacy-preserving alternatives rather than as accuracy-maximising classifiers.
The classical centralized baselines provide an important check that the leakage-controlled task is genuinely difficult.

Figure~\ref{fig:fl_convergence} visualises the convergence dynamics from the original FL runs; Table~\ref{tab:formal} should be treated as the primary leakage-controlled comparison.

\begin{figure}[t]
    \centering
    \includegraphics[width=\columnwidth]{figures/federated_results.png}
    \caption{Federated-learning diagnostics for formalization prediction, summarising convergence behaviour, aggregate performance, and client-level variation from the original runs.}
    \label{fig:fl_convergence}
\end{figure}

\begin{table}[t]
    \centering
    \caption{Leakage-controlled formalization prediction results (mean~$\pm$~std).}
    \label{tab:formal}
    \resizebox{\columnwidth}{!}{%
    \begin{tabular}{@{}lrrr@{}}
        \toprule
        Method & Accuracy & F1 & AUC \\
        \midrule
        Logistic Regression & .733$\pm$.033 & .833$\pm$.020 & .615$\pm$.097 \\
        Random Forest & .819$\pm$.024 & .898$\pm$.014 & .800$\pm$.026 \\
        XGBoost & .811$\pm$.027 & .892$\pm$.018 & .779$\pm$.042 \\
        Local MLP (non-FL) & .811$\pm$.036 & .888$\pm$.024 & .676$\pm$.076 \\
        Centralized MarketNet & .719$\pm$.023 & .816$\pm$.019 & .681$\pm$.042 \\
        Local-only MarketNet (avg) & .694$\pm$.019 & .800$\pm$.018 & .518$\pm$.031 \\
        FedAvg & .683$\pm$.058 & .792$\pm$.045 & .594$\pm$.042 \\
        FedProx ($\mu{=}0.1$) & .703$\pm$.063 & .806$\pm$.046 & .663$\pm$.044 \\
        \bottomrule
    \end{tabular}
    }
\end{table}

The augmentation sensitivity analysis uses SMOTE to expand the 358-row WBES formalization dataset to 1,764 synthetic training rows in exploratory runs.
Because these synthetic rows can preserve local decision boundaries from the original small sample, we do not use the SMOTE result as primary evidence.
The sensitivity analysis is instead interpreted as an upper-bound diagnostic for whether a richer, balanced field dataset could support stronger formalization prediction.

\begin{table}[t]
    \centering
    \caption{Augmentation sensitivity for formalization prediction.}
    \label{tab:augment}
    \resizebox{\columnwidth}{!}{%
    \begin{tabular}{@{}lccc@{}}
        \toprule
        Method & A: FeatSel & B: SMOTE & C: Leakage-controlled \\
        & (358$\times$30) & (1764$\times$30) & (358 original rows) \\
        \midrule
        Centralized & .719$\pm$.023 & .62$\pm$.06 & .719$\pm$.023 \\
        Local-only & .694$\pm$.019 & .59$\pm$.05 & .694$\pm$.019 \\
        FedAvg & .683$\pm$.058 & .57$\pm$.06 & .683$\pm$.058 \\
        FedProx & .703$\pm$.063 & .60$\pm$.05 & .703$\pm$.063 \\
        \bottomrule
    \end{tabular}
    }
\end{table}

\subsection{Federated Price Prediction}

Table~\ref{tab:price} reports price prediction MSE with proper temporal held-out evaluation.
FedAvg and FedProx show competitive test MSE relative to centralised training, while local-only models overfit severely.
This provides evidence for cross-market knowledge sharing: federated models generalise better to unseen future time periods than isolated local models in these experiments.

\begin{table}[t]
    \centering
    \caption{Price prediction test MSE across temporal resampling runs.}
    \label{tab:price}
    \begin{tabular}{@{}lrr@{}}
        \toprule
        Method & Rwanda & Uganda \\
        \midrule
        ARIMA & .74$\pm$.09 & .58$\pm$.07 \\
        Prophet & .69$\pm$.08 & .54$\pm$.06 \\
        Centralized MLP & .53$\pm$.05 & .41$\pm$.04 \\
        Local MLP & 1.71$\pm$.22 & .85$\pm$.11 \\
        FedAvg & .56$\pm$.06 & .42$\pm$.05 \\
        FedProx ($\mu{=}0.1$) & .55$\pm$.05 & .42$\pm$.04 \\
        \bottomrule
    \end{tabular}
\end{table}

In the cross-border comparison, Uganda's markets exhibit lower MSE despite fewer clients (203 vs.\ 426), attributable to lower commodity diversity and less volatile price dynamics.
Both countries show the same broad ordering: centralized and federated MLPs improve on local-only training and classical univariate baselines, indicating that federated aggregation can support generalisation under the modeled temporal split.
Figure~\ref{fig:benchmark} presents the convergence curves and per-market MSE breakdown, revealing that the hardest markets for both countries are those near refugee camps (e.g., Nyabiheke, Kyaka~II), which experience volatile demand shocks.

\begin{figure*}[t]
    \centering
    \includegraphics[width=\textwidth]{figures/comparative_benchmark.png}
    \caption{Cross-country price forecasting results showing that federated models remain competitive across Rwanda and Uganda, with larger errors in more volatile market settings.}
    \label{fig:benchmark}
\end{figure*}

\subsection{Economic and Spatial Analysis}

Table~\ref{tab:viability} quantifies vendor viability across the three market scenarios.
Scenario analysis suggests potential income gains under assumed occupancy and demand conditions, while reducing modeled displacement risk and providing digital visibility.
Modern markets yield \emph{negative} net income ($-720{,}000$~RWF/year) due to rent exceeding revenue capacity.

\begin{table}[t]
    \centering
    \caption{Scenario-based vendor viability estimates (annual, RWF).}
    \label{tab:viability}
    \begin{tabular}{@{}lrrrl@{}}
        \toprule
        Scenario & Revenue & Cost & Net & Risk \\
        \midrule
        Street hawking & 250K & 50K & 25K & 30\% disp. \\
        AI modular unit & 625K & 10K & 178K & 0\% \\
        Modern market & 1,600K & 1,200K & $-$720K & 0\% \\
        \bottomrule
    \end{tabular}
\end{table}

\begin{table}[t]
    \centering
    \caption{Income sensitivity to assumed foot-traffic multiplier.}
    \label{tab:sensitivity}
    \resizebox{\columnwidth}{!}{%
    \begin{tabular}{@{}lrrr@{}}
        \toprule
        Foot-traffic multiplier & Revenue & Net income & Gain vs. street \\
        \midrule
        1.5$\times$ & 375K & 102.5K & 4.1$\times$ \\
        2.0$\times$ & 500K & 140.0K & 5.6$\times$ \\
        2.5$\times$ & 625K & 177.5K & 7.1$\times$ \\
        \bottomrule
    \end{tabular}
    }
\end{table}

The national average inclusion gap is $\Gamma = 149{,}175$~RWF/year per vendor, with total annual loss from informality of 33,820~M~RWF (\$24.2M).
With digital demand-pull ($\delta = 0.85$), urban districts experience amplified gaps due to higher $\beta_d$, highlighting that the digital divide \emph{widens} the inclusion gap.
Figure~\ref{fig:economic} summarises these modeled economic effects across vendor income, aggregate deployment impact, and annual cost--benefit structure.
The modular-unit scenario performs better than baseline scenarios under the stated assumptions because it combines lower recurring costs with demand-pull from improved visibility and location.
The comparison should therefore be read as a scenario estimate for policy planning, not as evidence of field-validated causal income gains.

\begin{figure}[t]
    \centering
    \includegraphics[width=\columnwidth]{figures/economic_impact.png}
    \caption{Scenario-based economic estimates indicating that modular units can improve vendor viability under assumed demand and cost conditions, while modern market rents remain constraining.}
    \label{fig:economic}
\end{figure}

Under the pilot budget of 100 modular units (5,000 vendor slots), risk-weighted spatial allocation (Eq.~\ref{eq:allocation}) improves on population-proportional placement.
Kigali's three districts (Gasabo, Kicukiro, Nyarugenge) receive 32 units serving 1,600 vendors, reflecting their combined 57,440 informal enterprises.
National coverage reaches 2.2\% of informal vendors at pilot scale.
Figure~\ref{fig:spatial_optimization} illustrates how the allocation rule shifts units toward districts with high informal-enterprise density and lower modeled viability risk.
This produces a deployment pattern that is more targeted than a purely population-proportional allocation, while still exposing the limited reach of a 100-unit pilot relative to national demand.

\begin{figure}[t]
    \centering
    \includegraphics[width=\columnwidth]{figures/spatial_optimization.png}
    \caption{Risk-weighted spatial optimisation prioritises districts where modular units can reach many informal vendors while maintaining economic viability.}
    \label{fig:spatial_optimization}
\end{figure}

\subsection{Market Graph Analysis}

Rwanda's spatial market graph (320 nodes, 27,048 edges within 80\,km) reveals a dense, well-connected market network with average degree 169.1, while Uganda's sparser graph (96 nodes, 597 edges, degree 12.4) reflects its fewer WFP-monitored markets.
Mean catchment population per market is 345,155 (Rwanda) and 148,722 (Uganda).
Average digital readiness is 0.522 (Rwanda) and 0.546 (Uganda), suggesting comparable mobile money and internet penetration.
The mathematical framework integrates these spatial and economic factors into a unified policy tool.

%=================================================================
\section{Discussion}
\label{sec:discussion}
%=================================================================

Informal vendors are among the most privacy-sensitive economic actors, making federated learning relevant for market analytics. They avoid formal systems precisely because data sharing with authorities risks taxation, eviction, or harassment \cite{chen2023rethinking}.
FL enables collective intelligence without requiring vendors to surrender individual transaction data.
Our results indicate that this privacy preservation can preserve competitive price-forecasting performance, although formalization prediction becomes difficult once registration-derived leakage is removed.
This is consistent with recent FL surveys showing that well-tuned federated methods can match centralised training for tabular and time-series tasks \cite{kairouz2021advances,wen2023survey}.
Unlike many AI-for-good papers, this work addresses income generation for privacy-sensitive informal workers using deployable, low-resource infrastructure that does not require vendors to surrender personal data.

The modular unit proposition is evaluated as a scenario analysis rather than as a deployed intervention. Potential income gains from modular units versus street hawking are driven by two modeled mechanisms: (i)~a demand-pull foot-traffic multiplier explored at 1.5$\times$, 2.0$\times$, and 2.5$\times$, and (ii)~shared infrastructure costs (10,000~RWF/year per vendor vs.\ 50,000~RWF for street overhead or 1,200,000~RWF for modern market rent).
The digital demand-pull coefficient $\beta_d$ is therefore interpreted as a scenario parameter rather than an empirically validated effect size.

Several limitations shape the interpretation of these results. First, there is no live deployment yet; all FL clients are simulated partitions, not real distributed devices.
Second, economic gains are modeled under assumed conditions, not field-validated.
Third, SMOTE augmentation expands 358 records to 1,764; results may reflect synthetic manifold memorization rather than true generalization.
Fourth, WBES data covers only formal enterprises, so the smallest informal vendors are underrepresented.
Fifth, the pilot allocation of 100 units covers only 2.2\% of informal vendors, and scaling requires substantial investment.
Future work includes a physical pilot and field validation of the foot-traffic multiplier.

The risk-weighted allocation model also raises fairness considerations because it explicitly penalises deployments in low-viability districts, reducing low-return investment but potentially under-serving the most marginalised areas.
Corbett-Davies and Goel~\shortcite{corbettdavies2023measure} highlight that optimising for aggregate utility can entrench inequality; we mitigate this by ensuring a minimum allocation of one unit per province.
Dataset representativeness is also a concern \cite{scheuerman2021datasets}: the WBES covers only formal enterprises, and future work must incorporate informal survey data.

The work aligns with SDG~8 (decent work and economic growth), SDG~9 (industry, innovation, and infrastructure), SDG~10 (reduced inequalities), and SDG~11 (sustainable cities and communities) \cite{un2015sdg}.
Vinuesa et al.~\shortcite{vinuesa2020sdg} show that AI can be an enabler for 134 of 169 SDG targets; our framework offers a concrete application at the intersection of AI, urbanisation, and economic inclusion.

%=================================================================
\section{Conclusion}
\label{sec:conclusion}
%=================================================================

We presented a federated market intelligence framework integrated with AI-optimised modular retail units to address the inclusion gap facing 229,764 informal enterprises in Rwanda.
Our mathematical formalisation of vendor viability ($N$), inclusion gap ($\Gamma$), digital demand-pull ($\beta$), and risk-weighted allocation ($S$) provides a structured basis for market infrastructure investment decisions.

Federated learning indicates feasibility for multi-commodity price forecasting (FedProx test MSE: $0.55\pm0.05$ for Rwanda, $0.42\pm0.04$ for Uganda), while leakage-controlled formalization prediction highlights the limits of small, proxy-rich survey data.
Scenario analysis suggests potential income gains from modular units versus street hawking, and risk-weighted spatial placement offers a scalable alternative to the economically unviable modern market model under the stated assumptions.

Future work will: (i)~deploy a physical pilot with 10 modular units across Kigali districts; (ii)~integrate differential privacy \cite{wei2022dpfl} into the FL pipeline; (iii)~extend the price forecasting model with Transformer-based architectures \cite{zhou2021informer,wu2023timesnet} and graph neural networks \cite{wu2021graphsurvey} to exploit the spatial market graph; and (iv)~develop a real-time mobile application connecting the federated intelligence layer to vendor decision-making.

%=================================================================
\section*{Ethical Statement}
%=================================================================

This research uses publicly available aggregate statistics and anonymised survey data.
No personally identifiable vendor information was collected or processed.
The federated learning framework is designed explicitly to preserve data sovereignty of vulnerable informal workers.
We acknowledge that AI-driven market optimisation could inadvertently disadvantage vendors in low-digital-readiness areas; our risk-weighted allocation model explicitly mitigates this by down-weighting high-risk placements.

%=================================================================
\bibliographystyle{named}
\bibliography{references}

\end{document}
